\documentclass[12pt,a4paper]{article}
\usepackage[utf8]{inputenc}
\usepackage{amsmath}
\usepackage{amsfonts}
\usepackage{amssymb}
\usepackage[brazil]{babel}
\usepackage{indentfirst}
\usepackage{url}
\usepackage{float}
\usepackage{color}
%%%%%%%%%%Codigos para o JAVA%%%%%%%%%%%%%%%%%%%%%%%%%%%%%%%%
\definecolor{pblue}{rgb}{0.13,0.13,1}
\definecolor{pgreen}{rgb}{0,0.5,0}
\definecolor{pred}{rgb}{0.9,0,0}
\definecolor{pgrey}{rgb}{0.46,0.45,0.48}
\usepackage{listings}
\lstset{language=Java,
  showspaces=false,
  showtabs=false,
  breaklines=true,
  showstringspaces=false,
  breakatwhitespace=true,
  commentstyle=\color{pgreen},
  keywordstyle=\color{pblue},
  stringstyle=\color{pred},
  basicstyle=\ttfamily,
  moredelim=[il][\textcolor{pgrey}]{\$\$},
  moredelim=[is][\textcolor{pgrey}]{\%\%}{\%\%}
}
%%%%%%%%%%%%%%%%%%%%%%%%%%%%%%%%Fim codigo JAVA%%%%%%%%%%%%%%
%%%%%%%%%%%Codigo geral%%%%%%%%%%%%%%%%%%%%%%%%%%%%%%%%%%%%%%
\definecolor{mygreen}{rgb}{0,0.6,0}
\definecolor{mygray}{rgb}{0.5,0.5,0.5}
\definecolor{mymauve}{rgb}{0.58,0,0.82}
\lstset{ %
  backgroundcolor=\color{white},   % choose the background color; you must add \usepackage{color} or \usepackage{xcolor}; should come as last argument
  basicstyle=\footnotesize,        % the size of the fonts that are used for the code
  breakatwhitespace=false,         % sets if automatic breaks should only happen at whitespace
  breaklines=true,                 % sets automatic line breaking
  captionpos=b,                    % sets the caption-position to bottom
  commentstyle=\color{mygreen},    % comment style
  deletekeywords={...},            % if you want to delete keywords from the given language
  escapeinside={\%*}{*)},          % if you want to add LaTeX within your code
  extendedchars=true,              % lets you use non-ASCII characters; for 8-bits encodings only, does not work with UTF-8
  frame=single,                    % adds a frame around the code
  keepspaces=true,                 % keeps spaces in text, useful for keeping indentation of code (possibly needs columns=flexible)
  keywordstyle=\color{blue},       % keyword style
  language=Octave,                 % the language of the code
  morekeywords={*,...},            % if you want to add more keywords to the set
  numbers=left,                    % where to put the line-numbers; possible values are (none, left, right)
  numbersep=5pt,                   % how far the line-numbers are from the code
  numberstyle=\tiny\color{mygray}, % the style that is used for the line-numbers
  rulecolor=\color{black},         % if not set, the frame-color may be changed on line-breaks within not-black text (e.g. comments (green here))
  showspaces=false,                % show spaces everywhere adding particular underscores; it overrides 'showstringspaces'
  showstringspaces=false,          % underline spaces within strings only
  showtabs=false,                  % show tabs within strings adding particular underscores
  stepnumber=1,                    % the step between two line-numbers. If it's 1, each line will be numbered
  stringstyle=\color{mymauve},     % string literal style
  tabsize=2,                       % sets default tabsize to 2 spaces
  title=\lstname                   % show the filename of files included with \lstinputlisting; also try caption instead of title
}
%%%%%%%%%%%%%%%%%%%%%%%%%%%%%%%%Fim codigo geral%%%%%%%%%%%%%
\RequirePackage{graphicx}
\title{O que é AJAX? Como se programa?}
\author{Daniel Moreira Cardoso \and Davi Ildeu de Faria \and Fernando Maciel \and Igor Justino Rodrigues \and José Gilvan Jacinto Junior \and Paulo Henrique Rodrigues Araujo}
 
\usepackage[left=3cm,right=3cm,top=2cm,bottom=2cm]{geometry}
\begin{document}
\begin{titlepage}
\begin{center}
\begin{figure}[htb]
                
                \label{figura:LogoIF}
        
                \centering
                \includegraphics[width=6cm]{recursos/imagens/logo.png} 
\end{figure}
Instituto Federal Goiano - Campus Ceres\\
Bacharelado em Sistemas de Informação\\
Prof. Me. Ronneesley Moura Teles\\\vspace{0.5cm}
Daniel Moreira Cardoso \\
Davi Ildeu de Faria \\
Fernando Maciel \\
Igor Justino Rodrigues \\
José Gilvan Jacinto Junior \\
Paulo Henrique Rodrigues Araujo \\
\vspace{5.0cm}
\textit{\textbf{\Large{O que é AJAX? Como se programa? }}}\\\vspace{0.5cm}
\vspace{9.5cm}
\end{center}
\end{titlepage}
\tableofcontents
\newpage
\begin{center}
\textbf{\Large{O que é AJAX? Como se programa?}}\\\vspace{0.5cm}
\end{center}
\section{Introdução}

O Ajax é um conjunto de tecnologias de desenvolvimento web que quando reunidas permite o desenvolvimento de uma aplicação melhor elaborada, intuitiva, e de melhor desempenho. Ajax Asynchronous JavaScript and XML(JavaScript assíncrono e XML) foi gerado por Jesse James Garret e o nome foi escolhido justamente para facilitar a explicação do modelo de interação diferente entre o navegador e o servidor web.

Asynchronous(assíncrono) é um tipo de comunicação que não ocorre exatamente ao mesmo tempo, não-simultânea. Ajax é o ato de carregar e renderizar uma página usando scripts rodando pelo lado do cliente carregando os dados em plano de fundo sem que haja a necessidade de atualização da página. Ajax não é uma tecnologia, mas sim um conjunto de tecnologias.

\section{Como o AJAX trabalha?}

Enquanto uma aplicação web convencional busca as informações no servidor e retorna ao cliente o AJAX não necessita dessa ação agilizando o processo de carregamento de dados, com o AJAX  toda a lógica e processamento é passado para o usuário, quando o usuário faz uma requisição quem busca e trás essas informações é o JavaScript de forma assíncrona evitando o famoso “reload” na tela. O tratamento dos dados, seu formato e exibição fica toda por conta do script que foi carregado inicialmente quando se acessou a página.

O processo inicial de carregamento é mais lento que de uma aplicação comum, pois muitas informações são pré-carregadas. Mas depois, somente os dados são carregados, tornando assim o site mais rápido.

\section{Aplicação}
A partir de um CRUD feito em PHP, para realizar as interações, realizamos um exemplo para conectar ao banco de dados e inserir sem a necessidade de atualização da página
\lstinputlisting[language=SQL]{recursos/codigo/banco.sql}
\lstinputlisting[language=PHP]{recursos/codigo/CRUD.php}
E pela página HTML enviamos os dados:
\lstinputlisting[language=HTML]{recursos/codigo/requisicao_banco_dados.html}
 
Já no próximo exemplo, informamos o nome e é retornado um alert com as horas, minutos e segundos atuais:
\lstinputlisting[language=HTML]{recursos/codigo/requisicao_dados.html}
Alert em PHP:
\lstinputlisting[language=PHP]{recursos/codigo/dados.php}


\section{Manipulando a resposta do servidor}
Lembre-se que quando você estava enviando a requisição, você forneceu o nome de uma função JavaScript que foi projetada para lidar com a resposta.

\lstinputlisting[language=HTML]{recursos/codigo/funcao.html}

Vamos ver o que essa função deve fazer. Primeiro, a função precisa checar o estado da requisição. Se o estado da requisição tem o valor igual a "4", significa que a resposta do servidor foi recebida por completo e está tudo OK para continuar o processo.

\lstinputlisting[language=HTML]{recursos/codigo/funcao2.html}

A lista completa dos valores readyState é a seguinte:

\begin{itemize}
\item 0 (não inicializado);
\item 1 (carregando);
\item 2 (carregado);
\item 3 (interativo);
\item 4 (completo).
\end{itemize}

\textbf{* readyState é algo como "estado de prontidão", mostra qual é o status do processo que
está sendo executado e se está sendo executado.}

A próxima coisa a se checar é o código de resposta do servidor HTTP. No exemplo a seguir, nós tratamos do retorno bem sucedido ou mal sucedido da requisição HTTP do AJAX, verificando se o código de resposta for 200.

\lstinputlisting[language=HTML]{recursos/codigo/funcao3.html}

Os principais Códigos de Requisição são:
\begin{itemize}
\item 200: "OK"
\item 403: "Proibido"
\item 404: Página Não Encontrada"
\end{itemize}

Agora, após você ter checado o estado da requisição e o código de status do HTTP da resposta, caberá a você fazer o que quiser com os dados que o servidor lhe enviou. Você tem duas opções para acessar esses dados:
\begin{itemize}


\item httpRequest.responseText – retorna a resposta do servidor como uma string de
texto
\item httpRequest.responseXML – retorna a resposta do servidor como um objeto
XMLDocument no qual você poderá percorrer usando as funções DOM do JavaScript.
\end{itemize}
Note que os passos acima são válidos somente se você usou uma solicitação assíncrona (terceiro parâmetro de open() foi definido como true). Se você usou um pedido síncrono
você não precisa especificar uma função, você pode acessar o retorno de dados pelo servidor diretamente depois de chamar send(), porque o script irá parar e esperar pela
resposta do servidor.







\section{Vantagens e Desvantagens}
\subsection{Vantagens}

\begin{itemize}
\item Uma das principais vantagens do uso do Ajax é o desempenho do site. O Ajax permite buscar  dados no servidor e mostrar em apenas uma parte da página assim não se fazendo necessário recarregar toda a página;

\item Outra vantagem bastante interessante sobre o Ajax é a compatibilidade. Não importa se a aplicação é em Java, PHP, Python ou outra linguagem de programação. (Sem o uso de frameworks);

\item Uma vantagem para o usuário é a interface. Normalmente o uso do Ajax proporciona uma melhor navegabilidade e facilidade de uso para o usuário.


\end{itemize}

\subsection{Desvantagens}

\begin{itemize}
\item Há uma boa compatibilidade com o Ajax e as linguagens de programação, porém quando usamos frameworks para implementar o Ajax podemos perder essa compatibilidade pois as vezes um framework pode interferir no funcionamento do outro. Como por exemplo a função \$ no Jquery e Prototype.

\end{itemize}

\section{Conclusão}

Com o desenvolvimento da pesquisa e implementação dos códigos Ajax foi possível perceber a real necessidade da utilização de tecnologias como esta em um mercado cada vez mais exigente e competitivo como é o nosso. Ainda contribuindo para em ganhos em desempenho e facilidade de uso (para cliente e desenvolvedor).



\section{Referências Bibliográficas}

\noindent \textbf{[1]} {https://imasters.com.br/artigo/10224/ajax/vantagens-e-desvantagens-do-uso-de-ajax-aspectos-praticos?trace=1519021197\&source=single.}\\\vspace{0.2cm}

\noindent \textbf{[2]} {https://developer.mozilla.org/pt-BR/docs/AJAX/Getting\_Started}\\\vspace{0.2cm}

\noindent \textbf{[3]} {http://www.devfuria.com.br/php/como-funcionam-os-metodos-get-e-post/}\\\vspace{0.2cm}




\noindent \textbf{[4]} {http://www.devfuria.com.br/php/como-funcionam-os-metodos-get-e-post/}\\\vspace{0.2cm}



\end{document}
%Modelo de código para inserir figura
%\begin{figure}[h]
%\centering
%\includegraphics[width=15cm]{logo.png}
%\label{4}
%\caption{Fonte:http://...; Acesso em 06/11/2017}
%\end{figure}